%%%%%%%%%%%%%%%%%%%%%%%%%%%%%%%%%%%%%%%%%%%%%%%%%%%%%%%%%%%%%%%%%
% PHIẾU HỌC TẬP - HÀM SỐ MŨ VÀ MÔ HÌNH TĂNG TRƯỞNG
% Phiên bản 2.0: Tiếp cận từ Toán học
%%%%%%%%%%%%%%%%%%%%%%%%%%%%%%%%%%%%%%%%%%%%%%%%%%%%%%%%%%%%%%%%%
\documentclass[12pt,a4paper]{article}

\usepackage[utf8]{inputenc}
\usepackage[T1]{fontenc}
\usepackage[vietnamese]{babel}
\usepackage{amsmath,amssymb}
\usepackage{array}
\usepackage{booktabs}
\usepackage{tabularx}
\usepackage{colortbl}
\usepackage{graphicx}
\usepackage{newtxtext, newtxmath}
\usepackage[scaled=0.92]{helvet}
\usepackage{microtype}

\usepackage[svgnames]{xcolor}
\definecolor{primary}{RGB}{0,102,153}
\definecolor{mathblue}{RGB}{0,51,102}
\definecolor{accent}{RGB}{255,102,0}
\definecolor{success}{RGB}{34,139,34}
\definecolor{TableHeader}{gray}{0.9}

\usepackage[left=1.5cm,right=1.5cm,top=2cm,bottom=2cm,headheight=15pt]{geometry}
\usepackage{fancyhdr}
\pagestyle{fancy}
\fancyhf{}
\renewcommand{\headrule}{\color{primary}\hrule width\textwidth height 1pt}
\fancyhead[L]{\color{primary}\footnotesize\sffamily\textbf{PHIẾU HỌC TẬP - HÀM SỐ MŨ}}
\fancyhead[R]{\color{primary}\footnotesize\sffamily\textbf{TOÁN 11}}
\fancyfoot[C]{\sffamily\thepage}

\usepackage{enumitem}
\usepackage{tcolorbox}
\tcbuselibrary{skins, breakable}
\usepackage[colorlinks=true,linkcolor=primary,urlcolor=accent]{hyperref}

\newtcolorbox{phieuhoctap}[1]{
    enhanced, breakable,
    colback=white, colframe=mathblue,
    boxrule=1.5pt, arc=3mm,
    fonttitle=\bfseries\sffamily\Large,
    coltitle=white, colbacktitle=mathblue,
    title=#1
}

\newtcolorbox{mathbox}{
    enhanced, breakable,
    colback=mathblue!5, colframe=mathblue!80,
    boxrule=1pt, arc=2mm,
    fonttitle=\bfseries\sffamily,
    title=\color{white}KIẾN THỨC TOÁN HỌC,
    colbacktitle=mathblue
}

\newtcolorbox{appbox}{
    enhanced, breakable,
    colback=accent!5, colframe=accent!80,
    boxrule=1pt, arc=2mm,
    fonttitle=\bfseries\sffamily,
    title=\color{primary}ỨNG DỤNG THỰC TIỄN
}

\newcolumntype{C}{>{\centering\arraybackslash}X}
\newcolumntype{L}{>{\raggedright\arraybackslash}X}
\newcommand{\header}[1]{\cellcolor{TableHeader}\sffamily\bfseries#1}

\setlist[itemize,1]{label=\textbullet, leftmargin=*, nosep}
\setlist[itemize,2]{label=--, leftmargin=*, nosep}

\begin{document}

%==============================================================
% PHIẾU HỌC TẬP SỐ 1: ÔN TẬP HÀM SỐ MŨ
%==============================================================
\begin{phieuhoctap}{PHIẾU HỌC TẬP SỐ 1: ÔN TẬP HÀM SỐ MŨ}

\noindent\textbf{Họ và tên:} \underline{\hspace{6cm}} \hfill \textbf{Nhóm:} \underline{\hspace{2cm}}

\vspace{0.5cm}
\begin{mathbox}
\textbf{Định nghĩa:} Hàm số mũ cơ số $a$ ($a > 0$, $a \neq 1$) có dạng: $y = a^x$

\textbf{Số $e$:} $e = \lim_{n \to \infty}\left(1 + \frac{1}{n}\right)^n \approx 2.71828...$

\textbf{Công thức tăng trưởng liên tục:}
$$\boxed{y(t) = y_0 \cdot e^{rt}}$$
\begin{itemize}
\item $y_0$: Giá trị ban đầu (tại $t = 0$)
\item $r$: Tốc độ tăng trưởng ($r > 0$: tăng, $r < 0$: giảm)
\item $t$: Thời gian
\end{itemize}
\end{mathbox}

\vspace{0.5cm}
\textbf{\color{mathblue}BÀI 1: Tính chất hàm số mũ} (Điền vào chỗ trống)

\begin{enumerate}[label=\alph*)]
\item $a^0 = \underline{\hspace{2cm}}$
\item $a^1 = \underline{\hspace{2cm}}$
\item $a^{x+y} = a^x \cdot \underline{\hspace{2cm}}$
\item $a^{x-y} = \dfrac{a^x}{\underline{\hspace{2cm}}}$
\item $(a^x)^y = a^{\underline{\hspace{2cm}}}$
\item Đồ thị $y = a^x$ (với $a > 1$) luôn đi qua điểm $(\underline{\hspace{1cm}}, \underline{\hspace{1cm}})$
\end{enumerate}

\vspace{0.5cm}
\textbf{\color{mathblue}BÀI 2: Tính giá trị} (Dùng máy tính, làm tròn 2 chữ số)

\begin{tabularx}{\textwidth}{|L|C|C|C|C|}
\hline
$x$ & 0 & 1 & 2 & 3 \\
\hline
$e^x$ & \underline{\hspace{2cm}} & \underline{\hspace{2cm}} & \underline{\hspace{2cm}} & \underline{\hspace{2cm}} \\
\hline
\end{tabularx}

\vspace{0.5cm}
\textbf{\color{mathblue}BÀI 3: Công thức tăng trưởng}

Gửi ngân hàng 100 triệu VND, lãi suất 5\%/năm (lãi kép liên tục).

Công thức: $A = P \cdot e^{rt}$ với $P = 100$ triệu, $r = 0.05$

\begin{enumerate}[label=\alph*)]
\item Sau 1 năm: $A = 100 \cdot e^{0.05 \times 1} = 100 \cdot e^{\underline{\hspace{1cm}}} = \underline{\hspace{3cm}}$ triệu
\item Sau 5 năm: $A = 100 \cdot e^{0.05 \times 5} = 100 \cdot e^{\underline{\hspace{1cm}}} = \underline{\hspace{3cm}}$ triệu
\item Sau 10 năm: $A = 100 \cdot e^{0.05 \times 10} = 100 \cdot e^{\underline{\hspace{1cm}}} = \underline{\hspace{3cm}}$ triệu
\end{enumerate}

\vspace{0.5cm}
\textbf{\color{mathblue}BÀI 4: So sánh tăng trưởng}

\begin{tabularx}{\textwidth}{|C|C|C|C|}
\hline
\header{Thời gian (năm)} & \header{Tuyến tính $y = 100 + 5t$} & \header{Mũ $y = 100 \cdot e^{0.05t}$} & \header{Chênh lệch} \\
\hline
0 & 100 & 100 & 0 \\
\hline
10 & \underline{\hspace{2cm}} & \underline{\hspace{2cm}} & \underline{\hspace{2cm}} \\
\hline
20 & \underline{\hspace{2cm}} & \underline{\hspace{2cm}} & \underline{\hspace{2cm}} \\
\hline
\end{tabularx}

\textbf{Nhận xét:} \underline{\hspace{0.85\textwidth}}

\end{phieuhoctap}

\newpage
%==============================================================
% PHIẾU HỌC TẬP SỐ 2: ỨNG DỤNG VÀO DỊCH BỆNH
%==============================================================
\begin{phieuhoctap}{PHIẾU HỌC TẬP SỐ 2: ỨNG DỤNG VÀO MÔ HÌNH DỊCH BỆNH}

\noindent\textbf{Họ và tên:} \underline{\hspace{6cm}} \hfill \textbf{Nhóm:} \underline{\hspace{2cm}}

\vspace{0.5cm}
\begin{appbox}
\textbf{Tình huống:} Một thành phố có 10,000 dân. Xuất hiện 10 ca nhiễm bệnh X ban đầu.

\textbf{Áp dụng công thức Toán:} $I(t) = I_0 \cdot e^{rt}$
\begin{itemize}
\item $I_0 = 10$ (số ca ban đầu)
\item $r = \beta - \gamma = 0.3 - 0.1 = 0.2$ (tốc độ lây lan 20\%/ngày)
\item $\beta$: Tỷ lệ lây nhiễm, $\gamma$: Tỷ lệ hồi phục
\end{itemize}
\end{appbox}

\vspace{0.5cm}
\textbf{\color{accent}BÀI 1: Tính số ca nhiễm theo công thức mũ}

Công thức: $I(t) = 10 \cdot e^{0.2t}$

\begin{tabularx}{\textwidth}{|C|C|C|C|C|C|}
\hline
\header{Ngày $t$} & 0 & 5 & 10 & 15 & 20 \\
\hline
$I(t) = 10 \cdot e^{0.2t}$ & 10 & \underline{\hspace{1.5cm}} & \underline{\hspace{1.5cm}} & \underline{\hspace{1.5cm}} & \underline{\hspace{1.5cm}} \\
\hline
\end{tabularx}

\textbf{Nhận xét:} Số ca tăng rất nhanh theo thời gian vì \underline{\hspace{5cm}}

\vspace{0.5cm}
\textbf{\color{accent}BÀI 2: Sử dụng EpidemicSim}
\begin{itemize}
\item \textbf{Địa chỉ:} \href{https://stem-1h8.pages.dev/EpidemicSim.html}{stem-1h8.pages.dev/EpidemicSim.html}
\item Nhập: N = 10,000, $I_0$ = 10, $\beta$ = 0.3, $\gamma$ = 0.1
\item Nhấn ``Chạy mô phỏng''
\end{itemize}

\textbf{Ghi lại kết quả từ mô phỏng:}

\begin{tabularx}{\textwidth}{|C|C|C|C|C|C|}
\hline
\header{Ngày $t$} & 0 & 5 & 10 & 15 & 20 \\
\hline
Từ EpidemicSim & 10 & \underline{\hspace{1.5cm}} & \underline{\hspace{1.5cm}} & \underline{\hspace{1.5cm}} & \underline{\hspace{1.5cm}} \\
\hline
\end{tabularx}

\textbf{\color{accent}BÀI 3: So sánh công thức mũ với mô phỏng SIR}

\begin{enumerate}[label=\alph*)]
\item Giai đoạn đầu (ngày 0-10): Công thức mũ và mô phỏng có giống nhau không?

\underline{\hspace{\textwidth}}

\item Giai đoạn sau (ngày 15-20): Tại sao kết quả khác nhau?

\underline{\hspace{\textwidth}}

\item Đỉnh dịch (peak) xảy ra vào ngày thứ \underline{\hspace{2cm}} với \underline{\hspace{3cm}} ca
\end{enumerate}

\vspace{0.5cm}
\textbf{\color{accent}BÀI 4: Thử nghiệm can thiệp}

Thay đổi $\beta$ từ 0.3 xuống 0.2 (giãn cách xã hội) và quan sát:

\begin{tabularx}{\textwidth}{|L|C|C|}
\hline
& \header{Không can thiệp ($\beta=0.3$)} & \header{Giãn cách ($\beta=0.2$)} \\
\hline
Đỉnh dịch (số ca max) & \underline{\hspace{3cm}} & \underline{\hspace{3cm}} \\
\hline
Ngày đạt đỉnh & \underline{\hspace{3cm}} & \underline{\hspace{3cm}} \\
\hline
\end{tabularx}

\textbf{Kết luận:} Giãn cách xã hội giúp \underline{\hspace{6cm}}

\vspace{0.5cm}
\textbf{\color{primary}KẾT LUẬN TỔNG QUAN}

\textbf{Kiến thức Toán học:}
\begin{itemize}
\item Công thức tăng trưởng mũ: $y = y_0 \cdot e^{rt}$
\item Áp dụng để mô hình hóa bài toán thực tế
\end{itemize}

\textbf{Ứng dụng thực tiễn:}
\begin{itemize}
\item Lãi kép ngân hàng, tăng trưởng dân số, lây lan dịch bệnh
\item Công thức mũ đúng ở giai đoạn đầu, sau đó cần mô hình phức tạp hơn (SIR)
\end{itemize}

\end{phieuhoctap}

\newpage
%==============================================================
% RUBRIC ĐÁNH GIÁ
%==============================================================
\begin{center}
{\Large\bfseries\sffamily\color{primary} RUBRIC ĐÁNH GIÁ SẢN PHẨM}
\end{center}

\vspace{0.3cm}
\noindent\textbf{Tên nhóm:} \underline{\hspace{4cm}} \hfill \textbf{Người đánh giá:} \underline{\hspace{4cm}}

\vspace{0.5cm}
\begin{tabularx}{\textwidth}{|>{\raggedright\arraybackslash}p{2.5cm}|X|X|X|c|}
\hline
\rowcolor{TableHeader}
\sffamily\bfseries Tiêu chí & \sffamily\bfseries Xuất sắc (4đ) & \sffamily\bfseries Khá (3đ) & \sffamily\bfseries Cần cải thiện (1-2đ) & \sffamily\bfseries Điểm \\
\hline
\textbf{1. Kiến thức Toán} & Nắm vững hàm mũ $y=a^x$, tính đúng $e^{rt}$ & Hiểu cơ bản, tính đúng & Còn nhầm lẫn công thức & \underline{\hspace{1cm}} \\
\hline
\textbf{2. Mô hình hóa} & Viết đúng công thức $I=I_0 e^{rt}$ & Viết được với hỗ trợ & Chưa viết được & \underline{\hspace{1cm}} \\
\hline
\textbf{3. Sử dụng công nghệ} & Thành thạo mô phỏng, so sánh kịch bản & Sử dụng cơ bản & Cần hướng dẫn nhiều & \underline{\hspace{1cm}} \\
\hline
\textbf{4. Trình bày} & Mạch lạc, có biểu đồ minh họa & Rõ ràng, có số liệu & Rời rạc & \underline{\hspace{1cm}} \\
\hline
\end{tabularx}

\vspace{0.5cm}
\noindent\textbf{TỔNG ĐIỂM:} \underline{\hspace{2cm}} / 16 điểm

\vspace{0.5cm}
\noindent\textbf{Nhận xét:}

\underline{\hspace{\textwidth}}

\underline{\hspace{\textwidth}}

\end{document}
